%------------------------------------------------------------------------------%

\usepackage{integracao_e_series}
\usepackage{systeme}

\title{Integração por Frações Parciais}

%------------------------------------------------------------------------------%
\begin{document}

\addtitle

%------------------------------------------------------------------------------%
\section{Objetivo e Preliminares}

%------------------------------------------------------------------------------%
\begin{frame}{Objetivo}
  Integrar funções da forma
  \[
    f(x) = \frac{P(x)}{Q(x)}
  \]
  onde $P$ e $Q$ são polinômios
\end{frame}

%------------------------------------------------------------------------------%
\begin{frame}{Polinômios}
  Se o polinômio $P$, de grau $n$, possui $n$ raízes reais distintas $x_i$
  \[
    P(x) = a(x-x_1)(x-x_2)\cdots(x-x_n)
  \]
  \pause
  Se a raiz $x_k$ se repete $p$ vezes
  \[
    P(x) = a(x-x_1)(x-x_2)\cdots(x-x_k)^p
  \]
  \pause
  Se $P$ possui raízes complexas
  \[
    P(x) = a(x-x_1)(x-x_2)\cdots(x^2+bx+c)
  \]
\end{frame}

%------------------------------------------------------------------------------%
\begin{frame}{Integral Relevante -- 1}
  \begin{align*}
    \onslide<+->{\int \frac{1}{x+a}\, dx}
    \onslide<+->{& = \int \frac{1}{u}\, du  & u = x+a \quad du = dx \\[5mm]}
    \onslide<+->{& = \ln\abs{u} + C \\[5mm]}
    \onslide<+->{& = \ln\abs{x+a} + C}
  \end{align*}
\end{frame}

%------------------------------------------------------------------------------%
\begin{frame}{Integral Relevante -- 2}
  \begin{align*}
    \onslide<+->{\int \frac{1}{x^2+a^2}\, dx}
    \onslide<+->{& = \int \frac{1}{a^2\left(\left(\frac{x}{a}\right)^2 + 1\right)} \,dx \\[2mm]}
    \onslide<+->{& = \frac{1}{a} \int \frac{1}{\left(\frac{x}{a}\right)^2 + 1} \,\frac{dx}{a}}
    \onslide<+->{& u = \frac{x}{a} \qquad du = \frac{dx}{a}  \\[2mm]}
    \onslide<+->{& = \frac{1}{a} \int \frac{1}{u^2+1}\, du  \\[2mm]}
    \onslide<+->{& = \frac{1}{a} \arctan(u) + C    \\[2mm]}
    \onslide<+->{& = \frac{1}{a} \arctan\left(\frac{x}{a}\right) + C}
  \end{align*}
\end{frame}

%------------------------------------------------------------------------------%
\begin{frame}{Exemplo}
  Integrar \(\int \frac{x^3+x}{x-1}\,dx\)

  \pause
  \[
    f(x) = \frac{P(x)}{Q(x)}
    \qquad\qquad
    P(x) = x^3+x
    \qquad\qquad
    Q(x) = x-1
  \]

  \pause
  Como o grau de $P$ é maior do que o de $Q$, calculamos a divisão
  \[
    f(x) = x^2 + x + 2 + \frac{2}{x-1}
  \]
\end{frame}

%------------------------------------------------------------------------------%
\begin{frame}{Exemplo}
  \begin{align*}
    \onslide<+->{\int \frac{x^3+x}{x-1} \,dx}
    \onslide<+->{& = \int x^2 + x + 2 + \frac{2}{x-1} \, dx \\[5mm]}
    \onslide<+->{& = \frac{x^3}{3}
      + \frac{x^2}{2}
      + 2x
      + 2 \ln\abs{x-1}
      + C}
  \end{align*}
\end{frame}

%------------------------------------------------------------------------------%
\section{Integração por Frações Parciais}

%------------------------------------------------------------------------------%
\begin{frame}{Integração por Frações Parciais}
  Integrar \(f(x)=\frac{P(x)}{Q(x)}\) onde $P$ e $Q$ são polinômios
  \pause
  \begin{enumerate}[<+->]
    \setlength\itemsep{1em}
    \item Se o grau de $P$ for maior ou igual que o grau de $Q$, dividir $P$ por $Q$

    \item Fatorar $Q(x)$

    \item Expandir $f$ em frações parciais
  \end{enumerate}
\end{frame}

%------------------------------------------------------------------------------%
\begin{frame}{Caso 1}
  $Q(x)$ é o produto de $n$ fatores lineares distintos
  \[
    Q(x) =
    \left( a_1 x + b_1 \right)
    \left( a_2 x + b_2 \right) \cdots
    \left( a_n x + b_n \right)
  \]

  \pause
  Expansão
  \[
    \frac{P(x)}{Q(x)}
    = \frac{A_1}{ a_1 x + b_1 }
    + \frac{A_2}{ a_2 x + b_2 }
    + \cdots
    + \frac{A_m}{ a_n x + b_n }
  \]
\end{frame}

\begin{frame}{Exemplo}
  \begin{align*}
    \onslide<+->{\frac{x-1}{x(x-1)(x+1)}}
    \onslide<+->{& = \frac{{\color{blue}A}}{x}
                   + \frac{{\color{blue}B}}{x-1}
                   + \frac{{\color{blue}C}}{x+1} \\[2mm]}
    \onslide<+->{& = \frac{{\color{blue}A}(x-1)(x+1)
                   + {\color{blue}B}x(x+1)
                   + {\color{blue}C}x(x-1)}{x(x-1)(x+1)} \\[7mm]}
    \onslide<+->{x-1 & = {\color{blue}A}(x-1)(x+1)
                       + {\color{blue}B}x(x+1)
                       + {\color{blue}C}x(x-1) \\}
    \onslide<+->{    & = x^2 \left( {\color{blue}A} + {\color{blue}B} + {\color{blue}C} \right)
                       + x   \left({\color{blue}B}-{\color{blue}C}\right)
                       - {\color{blue}A}}
  \end{align*}
\end{frame}

\begin{frame}{Exemplo}
\begin{columns}
\column{0.5\linewidth}
  \systeme[ABC]{
    A + B + C = 0,
        B - C = 1,
    A         = 1
  }
\column{0.5\linewidth}
  \pause
  \[
    A = 1
  \]
  \pause
  \[
    B = C
  \]
  \pause
  \[
    A + B + C = 1 + 2B = 0
  \]
  \pause
  \[
    B = -\frac{1}{2}
  \]
  \pause
  \[
    C = -\frac{1}{2}
  \]
\end{columns}
\end{frame}

\begin{frame}{Exemplo}
\begin{align*}
  \onslide<+->{\frac{x-1}{x(x-1)(x+1)}
  & = \frac{A}{x}
    + \frac{B}{x-1}
    + \frac{C}{x+1} \\[2mm]}
  \onslide<+->{& = \frac{1}{x}
    -\frac{1}{2} \frac{1}{x-1}
    -\frac{1}{2} \frac{1}{x+1}}
\end{align*}
\onslide<+->{Portanto}
\begin{align*}
  \onslide<3->{\int \frac{x-1}{x(x-1)(x+1)} \,dx
  & = \int \frac{1}{x} \,dx
    - \frac{1}{2} \int \frac{1}{x-1} \,dx
    - \frac{1}{2} \int \frac{1}{x+1} \,dx \\[2mm]}
  \onslide<+->{& = \ln\abs{x}
    - \frac{1}{2} \ln\abs{x-1}
    - \frac{1}{2} \ln\abs{x+1}
    + K}
\end{align*}
\end{frame}

%------------------------------------------------------------------------------%
\begin{frame}{Caso 2}
  $Q(x)$ é o produto de $m$ fatores lineares e alguns repetidos

  \pause
  Suponha que o primeiro fator $(a_1 x + b_1)$ seja repetido $p$ vezes
  \[
    Q(x)
    = \left(a_1 x + b_1\right)^p
      \left(a_2 x + b_2\right) \cdots
      \left(a_q x + b_q\right)
  \]

  \pause
  Expansão
  \[
    \frac{P(x)}{Q(x)}
    = \frac{A_1}{(a_1x+b_1)}
    + \frac{A_2}{(a_1x+b_1)^2}
    + \cdots
    + \frac{A_p}{(a_1x+b_1)^p}
    + \cdots
    + \frac{A_q}{a_q x + b_q}
  \]
\end{frame}

\begin{frame}{Exemplo}
  \[
    \frac{x^3-x-2}{x^3(x-1)}
    = \frac{A}{x}
    + \frac{B}{x^2}
    + \frac{C}{x^3}
    + \frac{D}{x-1}
  \]

  \pause
  \vspace{\baselineskip}
  \[
    \int \frac{x^3-x-2}{x^3(x-1)} \,dx
    = A \int \frac{1}{x} \,dx
    + B \int \frac{1}{x^2} \,dx
    + C \int \frac{1}{x^3} \,dx
    + D \int \frac{1}{x-1} \,dx
  \]
\end{frame}

%------------------------------------------------------------------------------%
\begin{frame}{Caso 3}
  $Q(x)$ contém fatores quadráticos irredutíveis que não se repetem
  \[
    Q(x) = \cdots  \left( ax^2 + bx + c \right)
    \qquad
    \text{ com } b^2-4ac<0
  \]

  \pause
  Expansão
  \[
    \frac{P(x)}{Q(x)}
    = \cdots
    + \frac{Ax+B}{ax^2 + bx + c}
  \]
\end{frame}
%------------------------------------------------------------------------------%

\begin{frame}{Exemplo}
  \[
    \frac{x}{(x-2)(x^2+1)(x^2+4)}
    = \frac{A}{x-2}
    + \frac{Bx+C}{x^2+1}
    + \frac{Dx+E}{x^2+4}
  \]
\end{frame}

%------------------------------------------------------------------------------%
\begin{frame}{Caso 4}
  $Q(x)$ contém fatores quadráticos irredutíveis repetidos
  \[
    Q(x) = \cdots  \left( ax^2 + bx + c \right)^r
    \qquad
    \text{ com } b^2-4ac<0
  \]

  \pause
  Expansão
  \[
    \frac{P(x)}{Q(x)}
    = \cdots
    + \frac{A_1x+B_1}{ax^2+bx+c}
    + \frac{A_2x+B_2}{(ax^2+bx+c)^2}
    + \cdots
    + \frac{A_rx+B_r}{(ax^2+bx+c)^r}
  \]
\end{frame}

\begin{frame}{Exemplo}
\[
  \frac{2x}{(x-9)(x^2+4)^2}
  = \frac{A}{x-9}
  + \frac{Bx+C}{x^2+4}
  + \frac{Dx+E}{(x^2+4)^2}
\]
\end{frame}

% 1
%------------------------------------------------------------------------------%
\addtocounter{exemplo}{1}
\section{Exemplo \theexemplo}

%------------------------------------------------------------------------------%
\begin{frame}{Exemplo \theexemplo}
  \onslide<+->{Calcule \(\int \frac{x^2+2x-1}{2x^3+3x^2-2x}dx\)}

  \onslide<+->{O grau do numerador $P$ é menor que o grau do denominador $Q$}

  \onslide<+->{Fatorando $Q$}
  \begin{align*}
    \onslide<3->{Q(x)
                 & = 2x^3+3x^2-2x \\}
    \onslide<+->{& = x\left(2x^2+3x-2\right) \\}
    \onslide<+->{& = x\, 2\left(x- \frac{1}{2}\right)\left(x+2\right) \\}
    \onslide<+->{& = x\left(2x-1\right)\left(x+2\right)}
  \end{align*}
  \onslide<+->{Caso 1 -- Fatores lineares sem repetição}
\end{frame}

%------------------------------------------------------------------------------%
\begin{frame}{Exemplo \theexemplo}
  \onslide<+->{\[
    \frac{x^2+2x-1}{2x^3+3x^2-2x}
    = \frac{A}{x}
    + \frac{B}{2x-1}
    + \frac{C}{x+2}
  \]}
  \onslide<+->{Multiplicando por \(2x^3+3x^2-2x = x\left(2x-1\right)\left(x+2\right)\)}
  \begin{align*}
    \onslide<+->{x^2+2x-1
    & = A(2x-1)(x+2)+Bx(x+2)+ Cx(2x-1) \\}
    \onslide<+->{& = (2A+B +2C)x^2+ (3A+2B-C)x-2A}
  \end{align*}
  \onslide<+->{\systeme[ABC]{
     2A + B +2C =  1,
     3A +2B - C =  2,
    -2A         = -1
  }}
  \onslide<+->{\[
    A= \frac{1}{2}  \qquad
    B= \frac{1}{5}  \qquad
    C=-\frac{1}{10}
  \]}
\end{frame}

%------------------------------------------------------------------------------%
\begin{frame}{Exemplo \theexemplo}
  \begin{align*}
  \onslide<+->{\int \frac{x^2+2x-1}{2x^3+3x^2-2x} dx
  & = \int
        \frac{\sfrac{1}{2}}{x}
      + \frac{\sfrac{1}{5}}{2x-1}
      - \frac{\sfrac{1}{10}}{x+2} \,dx \\[5mm]}
  \onslide<+->{& = \frac{1}{2}  \int \frac{1}{x} \,dx
    + \frac{1}{5}  \int \frac{1}{2x-1} \,dx
    - \frac{1}{10} \int \frac{1}{x+2} \,dx \\[5mm]}
  \onslide<+->{& = \frac{1}{2}  \ln\abs{x}
    + \frac{1}{10} \ln\abs{2x-1}
    - \frac{1}{10} \ln\abs{x+2}
    + C}
\end{align*}
\end{frame}


% 2
%------------------------------------------------------------------------------%
\addtocounter{exemplo}{1}
\section{Exemplo \theexemplo}

%------------------------------------------------------------------------------%
\begin{frame}{Exemplo \theexemplo}
  Calcule \(\int \frac{4x}{x^3-x^2-x+1}dx\)

  \pause
  O grau do numerador $P$ é menor que o grau do denominador $Q$

  \pause
  Raízes de $Q$: $x=1$ com multiplicidade 2 e $x=-1$

  \pause
  Fatorando $Q$
  \[
    Q(x)
    = x^3 - x^2 - x - 2
    = (x-1)^2(x+1)
  \]
  \pause
  Caso 2: Fatores lineares com repetição
\end{frame}

%------------------------------------------------------------------------------%
\begin{frame}{Exemplo \theexemplo}
  \onslide<+->{\[
    \frac{4x}{x^3-x^2-x-2}
    = \frac{A}{x-1}
    + \frac{B}{(x-1)^2}
    + \frac{C}{x+1}
  \]}
  \onslide<+->{Multiplicando os dois lados por \(x^3-x^2-x-2=(x-1)^2(x+1)\)}
  \begin{align*}
    \onslide<+->{4x
    & = A(x-1)(x+1)+B(x+1)+C(x-1)^2 \\}
    \onslide<+->{& = (A+C)x^2+(B-2C)x+(-A+B+C)}
  \end{align*}
  \onslide<+->{\systeme[ABC]{
     A    + C = 0,
        B -2C = 4,
    -A +B + C = 0
  }}
  \onslide<+->{\[
    A =  1 \qquad
    B =  2 \qquad
    C = -1
  \]}
\end{frame}

%------------------------------------------------------------------------------%
\begin{frame}{Exemplo \theexemplo}
  \begin{align*}
    \onslide<+->{\int \frac{4x}{x^3-x^2-x-2} \,dx
    & = \int \frac{1}{x-1}+ \frac{2}{(x-1)^2}-\frac{1}{x+1} \, dx \\[2mm]}
    \onslide<+->{& = \int \frac{1}{x-1} \,dx
      + \int \frac{2}{(x-1)^2} \,dx
      - \int \frac{1}{x+1} \,dx \\[2mm]}
    \onslide<+->{& = \ln \abs{x-1}- \frac{2}{x-1}- \ln\abs{x+1}+C \\[2mm]}
    \onslide<+->{& = \ln \abs{\frac{x-1}{x+1}} - \frac{2}{x-1}+C}
  \end{align*}
\end{frame}

% 3
%------------------------------------------------------------------------------%
\addtocounter{exemplo}{1}
\section{Exemplo \theexemplo}

%------------------------------------------------------------------------------%
\begin{frame}{Exemplo \theexemplo}
  Encontre \(\int \frac{2x^2-x+4}{x^3+4x}\,dx\)

  \pause
  O grau do numerador $P$ é menor que o grau do denominador $Q$

  \pause
  Fatorando $Q$
  \[
    Q(x)
    = x^3 + 4x
    = x\left( x^2 + 4 \right)
  \]
  \pause
  Caso 3: Fator quadrático irredutível sem repetição
\end{frame}

%------------------------------------------------------------------------------%
\begin{frame}{Exemplo \theexemplo}
  \onslide<+->{\[
    \frac{ 2x^2 - x + 4}{ x^3 + 4x }
    = \frac{2x^2 - x + 4}{x\left( x^2 + 4 \right)}
    = \frac{A}{x} + \frac{Bx+C}{x^2+4}
  \]}
  \onslide<+->{Multiplicando por \(x^3 + 4x=x\left(x^2+4\right)\)}
  \begin{align*}
    \onslide<+->{2x^2 - x + 4
    & = A\left( x^2 + 4 \right) + \left( Bx + C \right) x \\}
    \onslide<+->{& = \left( A + B \right) x^2 + Cx + 4A}
  \end{align*}
  \onslide<+->{\systeme[ABC]{
    A+B=2,
    C=-1,
    4A=4
    }}
  \onslide<+->{\[
    A =  1 \qquad
    B =  1 \qquad
    C = -1
  \]}
\end{frame}

%------------------------------------------------------------------------------%
\begin{frame}{Exemplo \theexemplo}
  \begin{align*}
    \onslide<+->{\int \frac{2x^2-x+4}{x^3+4x}dx
    & = \int \frac{1}{x} + \frac{x-1}{x^2+4} \,dx \\[3mm]}
    \onslide<+->{& = \int \frac{1}{x}dx + \int \frac{x}{x^2+4} dx - \int \frac{1}{x^2+4} dx \\[3mm]}
    \onslide<+->{& = \ln\abs{x}
    + \frac{1}{2} \ln\abs{x^2+4}
    - \frac{1}{2} \arctan\left(\frac{x}{2}\right)
    + C}
  \end{align*}
  \onslide<+->{Onde usamos
  \[
    u=x^2+4 \qquad \text{para calcular} \qquad\int \frac{x}{x^2+4} dx
  \]}
  \onslide<+->{\[
    \int \frac{1}{x^2+a^2}dx
    = \frac{1}{a} \arctan\left(\frac{x}{a}\right)+K
  \]}
\end{frame}

% 4
%------------------------------------------------------------------------------%
\addtocounter{exemplo}{1}
\section{Exemplo \theexemplo}

%------------------------------------------------------------------------------%
\begin{frame}{Exemplo \theexemplo}
  Calcule \(\int \frac{\sin(x)}{\cos^2(x)-3\cos(x)}dx\)

  \pause
  \[
    u  =  \cos(x)    \qquad
    \pause
    du = -\sin(x) dx
    \pause
    \quad\Rightarrow\quad
    \sin(x)dx = -du
  \]

  \pause
  \[
    \int \frac{\sin(x)}{\cos^2(x)-3\cos(x)}dx
    = -\int \frac{1}{u^2-3u}du
    \pause
    = -\int \frac{1}{u(u-3)}du
  \]
\end{frame}

%------------------------------------------------------------------------------%
\begin{frame}{Exemplo \theexemplo}
  Calculando \(\int \frac{1}{u(u-3)}du\)

  \pause
  Por frações parciais -- caso 1
  \[
    \frac{1}{u(u-3)}
    = \frac{A}{u} + \frac{B}{u-3}
  \]
  \pause
  Multiplicando por $u(u-3)$
  \[
    1
    = A(u-3)+Bu
    \pause
    = (A+B)u-3A
  \]
  \pause
  Resolvendo o sistema
  \[
    A = -\frac{1}{3} \qquad
    B =  \frac{1}{3}
  \]
\end{frame}

%------------------------------------------------------------------------------%
\begin{frame}{Exemplo \theexemplo}
\begin{align*}
  \onslide<+->{\int \frac{1}{u(u-3)}du
  & = \int \frac{-\sfrac{1}{3}}{u} + \frac{\sfrac{1}{3}}{u-3} \,du \\[2mm]}
  \onslide<+->{& = -\frac{1}{3} \int \frac{1}{u}   \,du
      +\frac{1}{3} \int \frac{1}{u-3} \,du \\[2mm]}
  \onslide<+->{& = - \frac{1}{3} \ln\abs{u}
      + \frac{1}{3} \ln\abs{u-3} + C \\[2mm]}
  \onslide<+->{& = \frac{1}{3} \ln \abs{\frac{u-3}{u}} + C }
\end{align*}
\end{frame}

%------------------------------------------------------------------------------%
\begin{frame}{Exemplo \theexemplo}
  \begin{align*}
    \onslide<+->{\int \frac{\sin(x)}{\cos^2(x)-3\cos(x)}dx
    & = -\int \frac{1}{u(u-3)}du \\[2mm]}
    \onslide<+->{& = -\frac{1}{3} \ln \abs{\frac{u-3}{u}} + C \\[2mm]}
    \onslide<+->{& =  \frac{1}{3} \ln \abs{\frac{u}{u-3}} + C \\[2mm]}
    \onslide<+->{& =  \frac{1}{3} \ln \abs{\frac{\cos(x)}{\cos(x)-3}} + C}
  \end{align*}
\end{frame}

%------------------------------------------------------------------------------%
\begin{frame}
  \tableofcontents
\end{frame}

%------------------------------------------------------------------------------%
\end{document}
%------------------------------------------------------------------------------%